\section{Grafen oxidová membrána}\label{s:grafen_oxidova_membrana}
V rámci této Kapitoly budeme termínem grafen oxid (GO) označovat obecně jakoukoliv naoxidovanou formu uhlíku. Ačkoliv by v počátečních fázích syntézy, kdy se jedná o neexfoliovaný produkt spíše než o rozptýlené monoatomární vrstvy, bylo terminologicky přesnější hovořit o oxidu grafitu, pro popis metodiky přípravy zavedeme jednotné značení GO. 

Spotřeba GO na výrobu jedné membrány se v závislosti na požadované tloušťce pohybuje v řádu desítek miligramů. Vzhledem k tomu, že standardní laboratorní syntéza poskytuje výtěžek v jednotkách gramů, je běžnou praxí příprava zásobní disperze, která si při vhodném skladování zachovává stabilitu po dobu několika let. Z časového hlediska vyžaduje samotná oxidace \numrange{2}{3} dny, zatímco následná purifikace představuje proces v řádu týdnů. Doba filtrace jedné membrány se pak pohybuje v rozmezí \qtyrange{5}{30}{\hour}.

%%%%%%%%%%%%%%%%%%%%%%%%%%%%%%%%%%%%%%%%%%%%%%%%%% Syntéza grafen oxidu %%%%%%%%%%%%%%%%%%%%%%%%%%%%%%%%%%%%%%%%%%%%%%%%%%
\subsection{Syntéza grafen oxidu}\label{ss:synteza_grafen_oxidu}
Níže uvedený postup popisuje autorskou syntézu a čištění jedné šarže GO. V experimentech byly využity i membrány vyrobené z dalších dvou zásobních disperzí, blíže specifikovaných v Kapitole ?.

%%%%%%%%%%%%%%%%%%%%%%%%%%%%%%%%%%%%%%%%%%%%%%%%%% Chemikálie a zařízení %%%%%%%%%%%%%%%%%%%%%%%%%%%%%%%%%%%%%%%%%%%%%%%%%%
\subsubsection{Chemikálie a zařízení}\label{sss:chemikalie_a_zarizeni}
Z použitých chemikálií, viz Tabulka \ref{tab:tour_method_chemicals}, jsou klíčové vlastnosti použitého grafitu. Výrobcem udávaná čistota se vztahuje ke kovovým nečistotám, deklarovaná velikost částic je \qtyrange{2}{15}{\micro\meter}. Dále byl použit led z~deionizované vody z~rozvodů VŠCHT.

\begin{table}[ht]
    \centering
    \caption{Seznam chemikálií použitých pro syntézu grafen oxidu.}
    \label{tab:tour_method_chemicals}
    \begin{tabular}{lllc S[table-format=3.1] @{\,} l}
        \toprule
        \textbf{Chemikálie} & \textbf{Čistota/Koncentrace} & \textbf{Výrobce} & \textbf{CAS} & \multicolumn{2}{c}{\textbf{Množství}}\\
        \midrule
        Grafit & \qty{99.9995}{\%} & Alfa Aesar & 7782-42-5 & 7.5 & \unit{\gram}\\
        \ce{H2SO4} & \qty{96}{\%} & Lach:Ner & 7664-93-9 & 900 & \unit{\milli\liter}\\
        \ce{H3PO4} & \qty{85}{\%} & Penta & 7664-38-2 & 100 & \unit{\milli\liter}\\
        \ce{KMnO4} & \qty{99}{\%} & Lach:Ner & 7722-64-7 & 45 & \unit{\gram}\\
        \ce{H2O2} & \qty{30}{\%} & Lach:Ner & 7722-84-1 & 200 & \unit{\milli\liter}\\
        % \ce{Ba(NO3)2} & 99 \% & Chemapol & 10022-31-8 \\ % pro čištění
        \bottomrule
    \end{tabular}
\end{table}

Pro chlazení reakční směsi byl použit cirkulační chladící agregát Julabo F12 s~řídící jednotkou ponorného termostatu Julabo MA, pro míchání a ohřev reakční směsi byla použita magnetická míchačka s topením IKA color squid.

%%%%%%%%%%%%%%%%%%%%%%%%%%%%%%%%%%%%%%%%%%%%%%%%%% Tourova metoda %%%%%%%%%%%%%%%%%%%%%%%%%%%%%%%%%%%%%%%%%%%%%%%%%%
\subsubsection{Tourova metoda}\label{sss:tourova_metoda}
Do \qty{900}{\milli\liter} koncentrované \ce{H2SO4} bylo za stálého míchání postupně přilito \qty{100}{\milli\liter} \ce{H3PO4}. Čirá směs byla následně s použitím termostatu Julabo vychlazena na teplotu \qty{0}{\celsius}.

Do vychlazeného systému bylo přidáno \qty{7.5}{\gram} mikročástic grafitu. Černá neprůhledná suspenze byla po dobu \qty{15}{\minute} homogenizována na magnetické míchačce při \qty{1000}{rpm}. 
Do systému bylo v~průběhu \qty{30}{\minute} postupně přidáno \qty{45}{g} \ce{KMnO4}. V~důsledku tvorby \ce{Mn2O7} došlo k nárůstu teploty na \qtyrange{25}{30}{\celsius} a tmavě zelenému olejovitému zbarvení.

Reakční směs byla řízeně ohřáta na \qty{50}{\celsius}, kde byla ponechána za stálého míchání po dobu \qty{6}{\hour}, dalších \qty{12}{\hour} byla směs míchána bez aktivního ohřívání. Směs průběžně měnila barvu z tmavě zelené na tmavě hnědou v důsledku postupující oxidace grafitu a tedy redukce \ce{Mn2O7} na \ce{MnO2}.

Reakční směs, již o pokojové teplotě, byla postupně přelita přes zhruba \qty{500}{\gram} ledu. Pro redukci zbytků \ce{Mn^{VII}} a \ce{Mn^{IV}} iontů byl po kapkách přidáván \ce{H2O2} za vzniku kyslíkové pěny. Přidáno bylo přibližně \qty{200}{\milli\liter} do dosažení charakteristického žlutého zbarvení grafen oxidu a vymizení pěnění.

%%%%%%%%%%%%%%%%%%%%%%%%%%%%%%%%%%%%%%%%%%%%%%%%%% Purifikace %%%%%%%%%%%%%%%%%%%%%%%%%%%%%%%%%%%%%%%%%%%%%%%%%%
\subsubsection{Purifikace}\label{sss:purifikace}
% první šarže - cite Jankovský 2016, doi.org/10.1016/j.apmt.2016.06.001
Výstupem syntézy je silně kyselá směs GO, síranů, fosforečnanů a manganatých iontů. Pro vyčištění produktu bylo použito opakované sedimentace či centrifugace a dekantace. Kvalita čištění byla v jednotlivých krocích sledována nejprve orientačně zbarvením lakmusového papírku v supernatantu, později pak až do negativní reakce supernatantu s dusičnanem barnatým jakožto indikátoru přítomnosti síranů.

V první fázi reakční směs byla natolik kyselá, že v řádu desítek hodin sedimentovala pouhou gravitací, viz Obrázek \ref{img:dekantace}. V průběhu prvních dvou týdnů proběhlo celkem pět cyklů dekantace, výstupem bylo zhruba \qty{200}{\milli\liter} suspenze, která již samovolně nesedimentovala. Poslední supernatant vykazoval pH \numrange{2.5}{3}.

\begin{figure}[ht]
    \centering
    \begin{subfigure}[b]{0.35\textwidth}
        \centering
        \includegraphics[width=\textwidth]{dekantace.png}
        \caption{Samovolná sedimentace}\label{img:dekantace}
        % https://profilab24.com/en/laboratory/thermostats/julabo-refrigerating-circulating-thermostat-ma#&gid=1&pid=1
    \end{subfigure}
    \hfill
    \begin{subfigure}[b]{0.3\textwidth}
        \centering
        \includegraphics[width=\textwidth]{zkumavka-suspenze.png}
        \caption{Prvotní centrifugace}\label{img:centrifugace_zacatek}
        % https://shop.brand.de/en/centrifuge-tubes-pp-with-screw-cap-p856.html
    \end{subfigure}
    \hfill
    \begin{subfigure}[b]{0.3\textwidth}
        \centering
        \includegraphics[width=\textwidth]{zkumavka-koloid.png}
        \caption{Konečná centrifugace}\label{img:centrifugace_konec}
    \end{subfigure}
    \caption{Průběh purifikace GO}
\end{figure}

V druhé fázi probíhalo čištění na centrifuze skupiny Sofergoup, konkrétně typ MPW~352RH s~rotorem na \qtyproduct{6x50}{\milli\liter} při fixním úhlu \ang{30} o~maximální frekvenci otáčení \qty{10000}{\rpm}. Použity byly kónické centrifugační zkumavky VWR s maximální relativní odstředivou silou \english{RCF, Relative Centrifugal Force} \qty{12500}{\xg}, dále zkumavky Brand s RCF \qty{15000}{\xg}. Produkt byl rovnoměrně rozdělen do dvou várek po~šesti zkumavkách. Zatímco jedna polovina vzorku podstupovala centrifugaci, druhá polovina byla po dekantaci a doplnění vodou vystavena působení ultrazvukové lázně KRAINTEK~12, která podporovala nejen samotné čištění homogenizací obsahu zkumavky, ale i žádoucí exfoliaci produktu. 

Každá zkumavka prošla dvanácti cykly centrifugace. V prvních cyklech došlo k~jasnému oddělení produktu a supernatantu už po \qty{15}{\minute} při \qty{5000}{\rpm} (\qty{2711}{\xg}), intenzita centrifugace byla postupně navyšována až na \qty{10000}{\rpm} (\qty{10732}{\xg}) po dobu \qty{2}{\hour}. Vlivem silně kyselého prostředí docházelo, i při dodržení maximální povolené hodnoty RCF, k rychlé degradaci materiálu zkumavek tvorbou mikrotrhlin \english{crazing}, viz Obrázek \ref{img:centrifugace_zacatek}, jejich životnost tak byla omezena na pouhé \numrange{2}{3} cykly. V~průběhu čištění postupně produkt přecházel ze suspenze do koloidní formy, která zaujímala zvyšující se část objemu zkumavky, viz Obrázek \ref{img:centrifugace_konec}. Supernatant po poslední centrifugaci již byl negativní na přítomnost síranů.

Po lyofilizaci vyčištěného produktu bylo získáno \qty{9.94}{\gram} produktu ve formě jemného prášku. Po odebrání malého množství pro analytické účely byla z produktu připravena zásobní disperze o objemu \qty{1000}{\milli\liter}, uchovávána ve dvou \qty{500}{\milli\liter} širokohrdlých neprůsvitných PE lahvích uskladněných za pokojové teploty mimo dosah světla.

%%%%%%%%%%%%%%%%%%%%%%%%%%%%%%%%%%%%%%%%%%%%%%%%%% Výroba membrán %%%%%%%%%%%%%%%%%%%%%%%%%%%%%%%%%%%%%%%%%%%%%%%%%%
\subsection{Výroba membrán}\label{ss:vyroba_membran}
\subsubsection{Přetlaková filtrace}
K vlastní výrobě membrán byla použita přetlaková filtrační cela Millipore \english{Solvent Resistant Stirred Cell} o průměru \qty{47}{\milli\meter}. Zařízení sestává z borosilikátového skleněného válce, který slouží jako zásobník disperze, a nerezové základny a víka, které jsou fixovány trojicí nosných sloupků s křížovými maticemi pro rovnoměrné dotažení. Maximální pracovní přetlak aparatury je limitován pevností skleněného válce a je výrobcem stanoven na \qty{6}{\bar}. Těsnění aparatury je zajištěno v horní části originálním fluoroelastomerovým O-kroužkem a v dolní části vyseknutým plochým těsněním z~obdobného materiálu.

Přívod plynu do cely byl realizován osazenou \qty[parse-numbers = false]{1/4}{\inchp} NPT nástrčnou rychlospojkou pro hadicové rozvody o vnějším průměru \qty{6}{\milli\meter}, jako zdroj plynu sloužil bezolejový pístový kompresor Puma MC 6206. Výstupní tlak byl udržován na hodnotě \qty{5}{\bar} pomocí redukčního ventilu s manometrem, který je součástí výstupní sestavy kompresoru. Převod z G~\qty[parse-numbers = false]{1/4}{\inchp} výstupu kompresoru na použitou \qty{6}{\milli\meter} polyuretanovou hadičku byl realizován redukčním šroubením Swagelok.

Pro výrobu membrány bylo dle požadované tloušťky použito adekvátní množství zásobní disperze, typicky \qtyrange{2}{10}{\milli\liter}, toto množství bylo v úzké \qty{50}{\milli\liter} kádince doplněno deionizovanou vodou na objem \qtyrange{40}{50}{\milli\liter}. Před samotnou filtrací byl obsah kádinky vždy po dobu \qtyrange{15}{30}{\minute} kontinuálně deaglomerován sonikátorem Qsonica~Q55. V důsledku sonikace došlo k zahřátí kádinky na zhruba \qtyrange{40}{45}{\celsius}.

Jako filtry byly použity Hahnemühle nylonové membránové filtry o průměru \qty{47}{\milli\meter} s nominální velikostí pórů \qty{0.45}{\micro\meter} o tloušťce \qty{100}{\micro\meter}. Filtr byl vložen na~dno suché filtrační cely, na filtr bylo vloženo ploché těsnění s použitím pinzety tak, aby nedošlo k poškození filtru. Následně byl sestaven zbytek cely, vrchní díl cely byl namontován až po vypuštění kompresoru. Po vystředění válce a rovnoměrném dotažení šroubů byla připravená disperze kvantitativně převedena do cely plnícím otvorem, kompresor byl zapnut až po úplném uzavření otvoru. Vzhledem k velmi nízké filtrační rychlosti byl filtrát jímán do absorpční textilie, která umožňovala jeho průběžné odpařování, textilie zároveň sloužila k překrytí celé aparatury pro zamezení nežádoucí fotoredukce GO. 

Filtrace nebyla ukončena bezprostředně po vymizení makroskopického sloupce kapaliny, nýbrž až po \qtyrange{12}{48}{\hour} od okamžiku, kdy povrch membrány nabyl vizuálně homogenního, matného a rovnoměrně vlhkého vzhledu.

Po odtlakování a rozebrání aparatury zůstávala sestava filtru s nanesenou membránou a plochého těsnění zpravidla přichycena ke spodní hraně skleněného válce. Pro zamezení popraskání membrány byl v průběhu další manipulace filtr vždy pevně přitisknut k pevné rovné podložce. K sejmutí válce z těsnění postačilo lehké zapáčení pinzetou. Oddělení těsnění od filtru probíhalo postupně po celém obvodu po zhruba \qtyrange{1}{2}{\centi\meter} úsecích, přičemž okraje úseku byly prsty zafixovány k podložce. Pro minimalizaci tahového napětí a prevenci mechanického poškození membrány bylo těsnění odstraňováno výhradně smykovým pohybem v rovině těsnění. 

\subsubsection{Sušení a oddělení membrány}
Pro zamezení mechanické deformace byl filtr s~membránou na Petriho misce po obvodu zafixován plastovým kroužkem a shora zatížen. Takto zajištěný vzorek byl následně na dobu alespoň \qty{48}{\hour} umístěn do exsikátoru nad nasycený roztok \ce{KCl} o~stabilní 85\% relativní vlhkosti.

Pro úpravu geometrie vzorku před jeho instalací byl ocelovým průbojníkem vyseknut kruhový výřez, v závislosti na kvalitě okrajů konkrétní membrány o průměru \qtyrange{34}{36}{\milli\meter}, přičemž nevyužitý zbytek materiálu byl uchován pro případné další analýzy. K následné separaci membrány od podkladového filtru byl s výhodou využit samotný permeametr. Vzorek byl membránou dolů položen na otevřenou spodní část permeační cely tak, aby plně překrýval plochu porézního kamene. Přivedením vakua do cely při zavřeném obtoku došlo k rovnoměrnému a celoplošnému přisátí membrány k podkladu, u takto zafixované membrány pak bylo možné bezpečně skalpelem narušit okraj a pinzetou plynule sloupnout filtr, aniž by došlo k poškození membrány.

Za účelem dlouhodobého uchování a prevence nežádoucí fotoredukce grafen oxidu byly hotové membrány uzavřeny v LDPE sáčcích a skladovány bez přístupu světla za pokojové teploty.

\iffalse

\begin{figure}[ht]
    \centering
    \begin{subfigure}[b]{0.45\textwidth}
        \centering
        \includegraphics[width=\textwidth]{termostat.jpg}
        \caption{Julabo F12-MA}\label{img:julabo_termostat}
        % https://profilab24.com/en/laboratory/thermostats/julabo-refrigerating-circulating-thermostat-ma#&gid=1&pid=1
    \end{subfigure}
    \hfill
    \begin{subfigure}[b]{0.45\textwidth}
        \centering
        \includegraphics[width=\textwidth]{ika michacka.jpg}
        \caption{IKA color squid}\label{img:ika_michacka}
    \end{subfigure}
    \caption{Zařízení použitá pro syntézu GO}
\end{figure}



\begin{figure}[ht]
    \centering
    \begin{subfigure}[b]{0.3\textwidth}
        \centering
        \includegraphics[width=\textwidth]{centrifuga.png}
        \caption{MPW 352RH}\label{img:MPW centrifuga}
        % https://www.unimed.cz/multifunkcni-chlazena-vyhrivana-centrifuga-mpw-352rh
    \end{subfigure}
    \hfill
    \begin{subfigure}[b]{0.3\textwidth}
        \centering
        \includegraphics[width=\textwidth]{rotor.jpg}
        \caption{Rotor fixní \ang{30}}\label{img:rotor}
        % https://www.unimed.cz/uhlovy-rotor-6x50-ml-s-hermetickym-vickem-2
    \end{subfigure}
    \hfill
    \begin{subfigure}[b]{0.3\textwidth}
        \centering
        \includegraphics[width=\textwidth]{kraintek.jpg}
        \caption{KRAINTEK 12}\label{img:kraintek_ultrazvuk}
        % https://www.kraintek.cz/prumyslove-ultrazvukove-cisticky-s-mechanickym-ovladanim/prumyslova-ultrazvukova-cisticka-k-12im-nerezova-mechanicke-ovladani-d197
    \end{subfigure}
    \caption{Zařízení použitá pro purifikaci GO}
\end{figure}
\begin{figure}[ht]
    \centering
    \begin{subfigure}[b]{0.3\textwidth}
        \centering
        \includegraphics[width=\textwidth]{millipore.png}
        \caption{Filtrační cela Millipore}\label{img:millipore}
        % https://www.sigmaaldrich.com/CZ/cs/product/mm/xfuf04701
    \end{subfigure}
    \hfill
    \begin{subfigure}[b]{0.3\textwidth}
        \centering
        \includegraphics[width=\textwidth]{kompresor.jpg}
        \caption{Puma MC 6206}\label{img:kompresor}
        % https://obchod.svet-piva.cz/zbozi/kompresor-puma-1-2hp-mc0206
    \end{subfigure}
    \hfill
    \begin{subfigure}[b]{0.3\textwidth}
        \centering
        \includegraphics[width=\textwidth]{sonikator.png}
        \caption{QSonica Q55}\label{img:sonikator}
        % https://www.sonicator.com/products/q55-sonicator
    \end{subfigure}
    \caption{Zařízení použitá pro výrobu membrán}
\end{figure}

\fi